\chapter{cute 之 Hopper前言}

\begin{center}\small 作者：reed（知乎 @reed-84-49，专栏「CUDA高性能编程」）\quad 原文：\url{https://zhuanlan.zhihu.com/p/1948048832794427610}\end{center}

\section{背景与动机}

之前的“cute之[1][2][3][4][5][6][7][8][9]”系列文章我们重点介绍了cute中的核心概念和抽象，并且基于这些概念和抽象，按照自低向上的方式逐步构建起了利用cute实现GEMM计算的核心框架，这些内容主要是以Ampere架构为基础进行的介绍。2022年3月，NVidia在GTC大会上发布了Ampere的后继硬件架构-Hopper架构[10]，并于同年10月份发布了其针对AI计算市场的旗舰级产品H100。如今已是2025年11月，经过3年的时间，该架构下的芯片已然成为各大数据中心的算力输出的核心产品，其构筑了当今AI计算的最核心力量。相较于前序的Ampere架构，Hopper架构上增加了很多功能feature，其编程范式也发生了一系列变化，这些feature和新范式组合在一起共同组成了核心矩阵计算能力。刚好我也由于工作上的原因有机会接触到Hopper架构的产品，所以形成此系列文章，一方面作为我个人学习认知的总结和沉淀，另一方面也分享给各位奋斗、挣扎[11]在CUDA进阶之路上的同仁们，希望它能澄清Hopper上的编程范式和结构，希望这个系列能够帮助到你们。

\section{系列内容规划}

总体构思上大概分为五篇来介绍该内容：

\begin{itemize}
  \item cute之Hopper MBarrier
  \item cute之Hopper TMA
  \item cute之Hopper WGMMA
  \item cute之Hopper 硬件流水线
  \item cute之Hopper 高效GEMM实现
\end{itemize}
